% page 166 du fac-simile (image p-165 du scan)
% bloc 162.6 x 229.0 mm ; marge gauche 8.0 mm
\pgc{-12.700mm}{0.000mm}{
\l{\sou{fakiro}. Nomo (en India) di la asketi mendikera di omna sekti}
\l{\cel{3}(Mohamedisti, Civaisti, Vishnuisti) qui probas aquirar}
\l{\cel{3}santeso per kontemplo e su-mutili. - DEFIRS.}
\l{}
\l{\sou{faksimilo}. Reprodukturo exakta de skriburo, de desegnuro,}
\l{\cel{3}sive manuale, sive per imprimo, grabo, e c. - DEFIRS.}
\l{}
\l{\sou{fakto}. I. To quo existas reale. - II. (\sou{Yurocienco}) To quo}
\l{\cel{3}eventas en afero, to quo judiciesas. - III. To quo even-}
\l{\cel{3}tis. - DEFIR.}
\l{}
\l{\sou{faktitivo}.\cel{2}(\sou{gram.})\cel{3}Qua indikas ke la subjekto di la verbo}
\l{\cel{3}eventigas la ago.}
\l{}
\l{\sou{faktoro}.\cel{2}I. (\sou{matem}.) En la multipliko, la multiplikanto e la}
\l{\cel{3}multiplikendo. - II. (\sou{metaf}.) To (o : ta) quo (o = qua)}
\l{\cel{3}efikas por la rezultajo. - DEFIRS.}
\l{}
\l{\sou{faktorio}. Kontoro di agenti di kompanio komercala en lando}
\l{\cel{3}stranjera. - DEFIS.}
\l{}
\l{\sou{faktorialo}. (\sou{matem.}) (\sou{algebro)}. Produto di qua la faktori esas}
\l{\cel{3}segun progresiono geometriala. - DE.}
\l{}
\l{\sou{\textquotedbl{}faktotum\textquotedbl{}}\cel{2}La homo qua, serve di ulu, agas irgo por plezar}
\l{\cel{3}a lu. - DEF.}
\l{}
\l{\sou{fakt}uro. (\sou{komerco}). Noto quan la vendisto furnisas a la kom-}
\l{\cel{3}printo di la vari quin ita livras, sur qua indikesas la}
\l{\cel{3}preco di ta vari. - DEFIS.}
\l{}
\l{\sou{fakultato}. I. (\sou{filoz}.) Povo di ento libera pri uzar sua agiveso.}
\l{\cel{3}funcionala. - II. (\sou{universitato)}.\cel{2}Korpo de profesori komi-}
\l{\cel{3}sita por kurso-docar ed examenar, en omna fako di la domeno}
\l{\cel{3}\textquotedbl{}doco supra-grada\textquotedbl{}. - DEFIRS.}
\l{}
\l{\sou{fakultativa}. Quan onu darfas agar o ne agar. - DEFIS.}
\l{}
\l{\sou{falar}. (\sou{netrans}.)Irar de supre ad infre, pro la efiko dagra-}
\l{\cel{3}vito. - DE.}
\l{}
\l{\sou{falango}.\cel{2}I. (\sou{anat.}) Singla ek la mikra osti longa, ye qui kon-}
\l{\cel{3}sistas la fingri di la manui e di la pedi. - II. (\sou{milito}).}
\l{\cel{3}(\sou{che la Greki antiqua}). Korpo de infantriani. - DEFIRS.}
\l{}
\l{\sou{falangino}. (\sou{anat.}) Falango duesma di la fingri trifalanga.}
\l{}
\l{\sou{falansterio.}\cel{1}(\sou{sociologio}).En la sistemo sociala da la Franco}
\l{\cel{3}Fourier : habiteyo di singla grupo de cent familii.- DEFIRS.}
\l{}
\l{\sou{falbalo}. Bendo de stofo plisita, frunsita, e c., qua uzesas}
\l{\cel{3}por ornar jupo, manteleto, shultro-tuko, qua fixigesas sur}
\l{\cel{3}ici per lua parto supra. - DFI.}
}
